\documentclass[11pt]{article}
\usepackage[a4paper,margin=26mm]{geometry}
\usepackage{amsmath,amssymb,mathtools}
\usepackage{unicode-math}
\setmainfont{Latin Modern Roman}
\setmathfont{Latin Modern Math}
\usepackage{xurl,hyperref}
\hypersetup{hidelinks,pdftitle={Problem statement}}
\newcommand{\sref}[2]{\hyperlink{source:#1:#2}{[S#2]}}
\newcommand{\cataloguescope}[1]{#1}
\begin{document}
\section{Hyperbolic geometry — Question 10.2}\label{statement}

\subsection{Definitions and mathematical statement}

\paragraph{Shared notation.}$\mathbb N=\{1,2,\ldots\}$, and $\mathbb Z,\mathbb Q,\mathbb R,\mathbb C$ denote the integers, rationals, reals and complex numbers. Any other conventions are specified locally.


A surface lamination $\Lambda$ is a closed subset of $M$ locally homeomorphic to a product of a disk with a closed transverse set. It is essential when it has no sphere leaf or torus leaf bounding a solid torus, and its complementary regions are irreducible, have incompressible boundary, and have no compressing monogon. Such a monogon is a compressing disk with one boundary arc on a boundary leaf and the other along a transverse interval fiber. It is genuine if some complementary region is not an interval bundle. A sublamination is a closed union of leaves. A manifold is atoroidal if it has no incompressible embedded torus other than a boundary-parallel one. Write $\widetilde M$ for its universal cover, and put a tilde on a foliation or lamination for its lift. For a lifted hyperbolic leaf $\lambda\cong\mathbb H^2$, continuous extension means that its inclusion in $\widetilde M\cong\mathbb H^3$ extends continuously to $\lambda\cup S^1_\infty(\lambda)\to\mathbb H^3\cup S^2_\infty$. Universal circles and other foliation-boundary constructions carry actions of $\pi_1(M)$, as do geometric ideal boundaries.
\paragraph{Statement recorded in the supplied source.}
In the hyperbolic-manifold setting, does every lifted leaf of an essential lamination have the continuous extension property? More generally, determine the relationship between group actions on ideal boundaries obtained from the foliated or laminated structure, including universal circles, and those on the geometric ideal boundaries of $\widetilde M$. Include the compatibility of these boundary constructions, not only existence of a leafwise extension. \sref{statement}{2}
\subsection{Short English statement}
Do essential-lamination leaves extend continuously to the ambient ideal boundary, and how do foliation boundaries and their group actions relate to geometric ones?
\subsection{Sources}
\begin{description}
\item[\hypertarget{source:statement:1}{S1}] ULAM.AI and the UnsolvedMath Contributors, \emph{UnsolvedMath: A Curated Collection of Open Mathematics Problems}, catalogue entry. CC BY 4.0. \url{https://huggingface.co/datasets/ulamai/UnsolvedMath}.
\item[\hypertarget{source:statement:2}{S2}] Danny Calegari, \emph{Problems in foliations and laminations of 3-manifolds}, version 0.78 (2002), Question 10.2; Definitions 1.1--1.10 and local definitions and remarks. \url{https://arxiv.org/abs/math/0209081v1}.
\end{description}
\label{end:statement}
\end{document}
